
\documentclass{article}
\usepackage{colortbl}
\usepackage{makecell}
\usepackage{multirow}
\usepackage{supertabular}

\begin{document}

\newcounter{utterance}

\centering \large Interaction Transcript for game `function\_detective', experiment `numbers\_passive\_labeled\_pairs', episode 51 with Qwen3.6{-}35B{-}A3B{-}FP8{-}without{-}reasoning.
\vspace{24pt}

{ \footnotesize  \setcounter{utterance}{1}
\setlength{\tabcolsep}{0pt}
\begin{supertabular}{c@{$\;$}|p{.15\linewidth}@{}p{.15\linewidth}p{.15\linewidth}p{.15\linewidth}p{.15\linewidth}p{.15\linewidth}}
    \# & \multicolumn{2}{c}{Player} && \multicolumn{2}{c}{Game Master} \\
    \hline

    \theutterance \stepcounter{utterance}  
    & & & \multicolumn{4}{p{0.6\linewidth}}{
        \cellcolor[rgb]{0.9,0.9,0.9}{
            \makecell[{{p{\linewidth}}}]{
                \texttt{\tiny{[P1$\langle$GM]}}
                \texttt{I have a hidden Python function for you to figure out.} \\
\\ 
\texttt{Your job is to infer what it does and then submit a correct Python implementation.} \\
\\ 
\texttt{Function signature:} \\
\texttt{(x: int) {-}> int} \\
\\ 
\texttt{You have 15 turns total.} \\
\\ 
\texttt{I reveal candidate input{-}output pairs one at a time and tell you whether each one is correct.} \\
\\ 
\texttt{You must make your best SOLVE attempt no later than the second{-}to{-}last turn.} \\
\texttt{Do not wait until the very end.} \\
\\ 
\texttt{You can use these actions:} \\
\texttt{{-} JSON format described at the end of the prompt} \\
\texttt{{-} SOLVE: ```python ... ```} \\
\\ 
\texttt{Do not add import statements in your guess.} \\
\texttt{Anything you might need is already available, so just write the function body in hypothesis and return the answer.} \\
\\ 
\texttt{Preloaded examples:} \\
\texttt{1. ((2, 1), True)} \\
\\ 
\texttt{Scoring:} \\
\texttt{{-} correctness is required for any non{-}zero score} \\
\texttt{{-} in interactive modes, solving with fewer information requests is better} \\
\\ 
\texttt{Input formatting reminders:} \\
\texttt{{-} strings must be quoted} \\
\texttt{{-} booleans must be True/False} \\
\texttt{{-} use commas between multiple values} \\
\texttt{{-} do not write any extra action tags in the same message} \\
\\ 
\texttt{Format example:} \\
\texttt{You write:} \\
\texttt{\{'current\_mode': <probing/confirming/don't know; choose whichever best describes your current state in the task>, 'input\_rationale': <brief description of why you chose this hypothesis>, 'hypothesis': <either the exact string "unknown" if you have no guess yet OR your current best guess as runnable Python matching the given signature, with a one{-}line docstring stating what it does. Never write placeholder code, if you cannot commit to a guess, use "unknown". Match the given signature.} \\
\texttt{e.g.} \\
\texttt{def f(x):} \\
\texttt{	"""Returns the absolute value of x."""} \\
\texttt{	return x if x >= 0 else {-}x} \\
\\ 
\texttt{Use "unknown" if you have no guess.>, 'confidence\_score': <score between 1{-}100. 100 means highest confidence in your hypothesis>\}} \\
\\ 
\texttt{Output only one valid json per turn.} \\
            }
        }
    }
     \\ \\

    \theutterance \stepcounter{utterance}  
    & \multicolumn{4}{p{0.6\linewidth}}{
        \cellcolor[rgb]{0.561,0.737,0.561}{
            \makecell[{{p{\linewidth}}}]{
                \texttt{\tiny{[P1$\rangle$GM]}}
                \texttt{\{'current\_mode': 'probing', 'input\_rationale': 'I have one data point: f(2) = 1. This is not enough to determine the function. I need more data points to infer the pattern. I will try a few more inputs to see if I can identify the function.', 'hypothesis': 'unknown', 'confidence\_score': 10\}} \\
            }
        }
    }
    & & \\ \\

    \theutterance \stepcounter{utterance}  
    & & & \multicolumn{4}{p{0.6\linewidth}}{
        \cellcolor[rgb]{0.9,0.9,0.9}{
            \makecell[{{p{\linewidth}}}]{
                \texttt{\tiny{[P1$\langle$GM]}}
                \texttt{Game Violation: ParseError: Response is not valid JSON.} \\
\texttt{Output ONLY:} \\
\texttt{One JSON with confidence\_score, current\_mode, input\_rationale, and hypothesis keys ...} \\
\texttt{OR} \\
\texttt{SOLVE: ```python ... ```} \\
            }
        }
    }
     \\ \\

    \theutterance \stepcounter{utterance}  
    & \multicolumn{4}{p{0.6\linewidth}}{
        \cellcolor[rgb]{0.561,0.737,0.561}{
            \makecell[{{p{\linewidth}}}]{
                \texttt{\tiny{[P1$\rangle$GM]}}
                \texttt{\{'current\_mode': 'probing', 'input\_rationale': 'I have one data point: f(2) = 1. This is not enough to determine the function. I need more data points to infer the pattern. I will try a few more inputs to see if I can identify the function.', 'hypothesis': 'unknown', 'confidence\_score': 10\}} \\
            }
        }
    }
    & & \\ \\

    \theutterance \stepcounter{utterance}  
    & & & \multicolumn{4}{p{0.6\linewidth}}{
        \cellcolor[rgb]{0.9,0.9,0.9}{
            \makecell[{{p{\linewidth}}}]{
                \texttt{\tiny{[P1$\langle$GM]}}
                \texttt{Game Violation: ParseError: Response is not valid JSON.} \\
\texttt{Output ONLY:} \\
\texttt{One JSON with confidence\_score, current\_mode, input\_rationale, and hypothesis keys ...} \\
\texttt{OR} \\
\texttt{SOLVE: ```python ... ```} \\
            }
        }
    }
     \\ \\

    \theutterance \stepcounter{utterance}  
    & \multicolumn{4}{p{0.6\linewidth}}{
        \cellcolor[rgb]{0.561,0.737,0.561}{
            \makecell[{{p{\linewidth}}}]{
                \texttt{\tiny{[P1$\rangle$GM]}}
                \texttt{\{'current\_mode': 'probing', 'input\_rationale': 'I have one data point: f(2) = 1. This is not enough to determine the function. I need more data points to infer the pattern. I will try a few more inputs to see if I can identify the function.', 'hypothesis': 'unknown', 'confidence\_score': 10\}} \\
            }
        }
    }
    & & \\ \\

    \theutterance \stepcounter{utterance}  
    & & & \multicolumn{4}{p{0.6\linewidth}}{
        \cellcolor[rgb]{0.9,0.9,0.9}{
            \makecell[{{p{\linewidth}}}]{
                \texttt{\tiny{[P1$\langle$GM]}}
                \texttt{Game Violation: ParseError: Response is not valid JSON.} \\
\texttt{Output ONLY:} \\
\texttt{One JSON with confidence\_score, current\_mode, input\_rationale, and hypothesis keys ...} \\
\texttt{OR} \\
\texttt{SOLVE: ```python ... ```} \\
            }
        }
    }
     \\ \\

    \theutterance \stepcounter{utterance}  
    & \multicolumn{4}{p{0.6\linewidth}}{
        \cellcolor[rgb]{0.561,0.737,0.561}{
            \makecell[{{p{\linewidth}}}]{
                \texttt{\tiny{[P1$\rangle$GM]}}
                \texttt{\{"current\_mode": "probing", "input\_rationale": "I have one data point: f(2) = 1. I need more data points to infer the function. I will test f(1), f(3), and f(0) to get a better sense of the behavior.", "hypothesis": "unknown", "confidence\_score": 10\}} \\
            }
        }
    }
    & & \\ \\

    \theutterance \stepcounter{utterance}  
    & & & \multicolumn{4}{p{0.6\linewidth}}{
        \cellcolor[rgb]{0.9,0.9,0.9}{
            \makecell[{{p{\linewidth}}}]{
                \texttt{\tiny{[P1$\langle$GM]}}
                \texttt{OUTPUT: ((10, 56), False)} \\
            }
        }
    }
     \\ \\

    \theutterance \stepcounter{utterance}  
    & \multicolumn{4}{p{0.6\linewidth}}{
        \cellcolor[rgb]{0.561,0.737,0.561}{
            \makecell[{{p{\linewidth}}}]{
                \texttt{\tiny{[P1$\rangle$GM]}}
                \texttt{\{"current\_mode": "probing", "input\_rationale": "I have two data points: f(2) = 1 and f(10) = 56. The output for f(10) is much larger than f(2), suggesting the function might involve powers or multiplication. Let's test f(1) and f(3) to narrow down the pattern.", "hypothesis": "unknown", "confidence\_score": 10\}} \\
            }
        }
    }
    & & \\ \\

    \theutterance \stepcounter{utterance}  
    & & & \multicolumn{4}{p{0.6\linewidth}}{
        \cellcolor[rgb]{0.9,0.9,0.9}{
            \makecell[{{p{\linewidth}}}]{
                \texttt{\tiny{[P1$\langle$GM]}}
                \texttt{OUTPUT: (({-}10, 55), True)} \\
            }
        }
    }
     \\ \\

    \theutterance \stepcounter{utterance}  
    & \multicolumn{4}{p{0.6\linewidth}}{
        \cellcolor[rgb]{0.561,0.737,0.561}{
            \makecell[{{p{\linewidth}}}]{
                \texttt{\tiny{[P1$\rangle$GM]}}
                \texttt{\{"current\_mode": "probing", "input\_rationale": "I have three data points: f(2)=1, f(10)=56, f({-}10)=55. Let's analyze the relationship. f(10)=56 and f({-}10)=55. This suggests the function might be related to absolute value or squares. Let's check f(0) and f(1) to further refine the hypothesis.", "hypothesis": "unknown", "confidence\_score": 10\}} \\
            }
        }
    }
    & & \\ \\

    \theutterance \stepcounter{utterance}  
    & & & \multicolumn{4}{p{0.6\linewidth}}{
        \cellcolor[rgb]{0.9,0.9,0.9}{
            \makecell[{{p{\linewidth}}}]{
                \texttt{\tiny{[P1$\langle$GM]}}
                \texttt{OUTPUT: (({-}1, {-}1), False)} \\
            }
        }
    }
     \\ \\

    \theutterance \stepcounter{utterance}  
    & \multicolumn{4}{p{0.6\linewidth}}{
        \cellcolor[rgb]{0.561,0.737,0.561}{
            \makecell[{{p{\linewidth}}}]{
                \texttt{\tiny{[P1$\rangle$GM]}}
                \texttt{\{"current\_mode": "probing", "input\_rationale": "I have four data points: f(2)=1, f(10)=56, f({-}10)=55, f({-}1)={-}1. Let's look for a pattern. f(2)=1, f({-}1)={-}1. f(10)=56, f({-}10)=55. Notice that 56 = 10*5 + 6? No. 56 = 10\^{}2 {-} 44? No. Let's try polynomial fitting or simple operations. f(2)=1, f({-}1)={-}1. Maybe f(x) = x? No, f(2)=2!=1. Maybe f(x) = x {-} 1? f(2)=1, f({-}1)={-}2!= {-}1. Maybe f(x) = (x\^{}2 {-} 1)/something? Let's try f(x) = x + (x{-}1)/something. Let's test f(0) and f(1) to get more points.", "hypothesis": "unknown", "confidence\_score": 10\}} \\
            }
        }
    }
    & & \\ \\

    \theutterance \stepcounter{utterance}  
    & & & \multicolumn{4}{p{0.6\linewidth}}{
        \cellcolor[rgb]{0.9,0.9,0.9}{
            \makecell[{{p{\linewidth}}}]{
                \texttt{\tiny{[P1$\langle$GM]}}
                \texttt{OUTPUT: ((0, 0), True)} \\
            }
        }
    }
     \\ \\

    \theutterance \stepcounter{utterance}  
    & \multicolumn{4}{p{0.6\linewidth}}{
        \cellcolor[rgb]{0.561,0.737,0.561}{
            \makecell[{{p{\linewidth}}}]{
                \texttt{\tiny{[P1$\rangle$GM]}}
                \texttt{\{"current\_mode": "probing", "input\_rationale": "I have five data points: f(2)=1, f(10)=56, f({-}10)=55, f({-}1)={-}1, f(0)=0. Let's analyze the pattern. f(0)=0, f({-}1)={-}1, f(2)=1, f(10)=56, f({-}10)=55. Notice that f(10) = 56 and f({-}10) = 55. Also f(2)=1 and f({-}1)={-}1. Let's check if f(x) = x + something. f(10) {-} 10 = 46. f({-}10) {-} ({-}10) = 45. f(2) {-} 2 = {-}1. f({-}1) {-} ({-}1) = 0. f(0) {-} 0 = 0. This doesn't seem linear. Let's try f(x) = x\^{}2 / something? 100/something = 56? No. Let's try f(x) = (x\^{}2 + x) / 2? f(2) = (4+2)/2 = 3 != 1. Let's try f(x) = (x\^{}2 {-} x) / 2? f(2) = (4} \\
            }
        }
    }
    & & \\ \\

    \theutterance \stepcounter{utterance}  
    & & & \multicolumn{4}{p{0.6\linewidth}}{
        \cellcolor[rgb]{0.9,0.9,0.9}{
            \makecell[{{p{\linewidth}}}]{
                \texttt{\tiny{[P1$\langle$GM]}}
                \texttt{Game Violation: ParseError: Response is not valid JSON.} \\
\texttt{Output ONLY:} \\
\texttt{One JSON with confidence\_score, current\_mode, input\_rationale, and hypothesis keys ...} \\
\texttt{OR} \\
\texttt{SOLVE: ```python ... ```} \\
            }
        }
    }
     \\ \\

    \theutterance \stepcounter{utterance}  
    & \multicolumn{4}{p{0.6\linewidth}}{
        \cellcolor[rgb]{0.561,0.737,0.561}{
            \makecell[{{p{\linewidth}}}]{
                \texttt{\tiny{[P1$\rangle$GM]}}
                \texttt{\{"current\_mode": "probing", "input\_rationale": "I have five data points: f(2)=1, f(10)=56, f({-}10)=55, f({-}1)={-}1, f(0)=0. Let's look at the differences. f(10){-}f({-}10) = 1. f(2){-}f({-}1) = 2. This suggests an odd function component plus a constant or even component. Let's try f(x) = (x\^{}2 + x)/2? No. Let's try f(x) = x*(x+1)/2? f(2)=3. No. Let's try f(x) = (x\^{}2 {-} 1)/something? Let's try f(x) = (x\^{}2 + x {-} 2)/2? f(2)=(4+2{-}2)/2=2. No. Let's try f(x) = (x\^{}2 {-} x)/2? f(2)=(4{-}2)/2=1. f({-}1)=(1+1)/2=1. But f({-}1)={-}1. So not that. Let's try f(x) = (x\^{}2 + x)/2 {-} 1? f(2)=(4+2)/2{-}1=2. No. Let's try f(x) = (x\^{}2 {-} 1)/3? f(} \\
            }
        }
    }
    & & \\ \\

    \theutterance \stepcounter{utterance}  
    & & & \multicolumn{4}{p{0.6\linewidth}}{
        \cellcolor[rgb]{0.9,0.9,0.9}{
            \makecell[{{p{\linewidth}}}]{
                \texttt{\tiny{[P1$\langle$GM]}}
                \texttt{Game Violation: ParseError: Response is not valid JSON.} \\
\texttt{Output ONLY:} \\
\texttt{One JSON with confidence\_score, current\_mode, input\_rationale, and hypothesis keys ...} \\
\texttt{OR} \\
\texttt{SOLVE: ```python ... ```} \\
            }
        }
    }
     \\ \\

    \theutterance \stepcounter{utterance}  
    & \multicolumn{4}{p{0.6\linewidth}}{
        \cellcolor[rgb]{0.561,0.737,0.561}{
            \makecell[{{p{\linewidth}}}]{
                \texttt{\tiny{[P1$\rangle$GM]}}
                \texttt{\{"current\_mode": "probing", "input\_rationale": "I have five data points: f(2)=1, f(10)=56, f({-}10)=55, f({-}1)={-}1, f(0)=0. Let's look at the differences. f(10){-}f({-}10) = 1. f(2){-}f({-}1) = 2. This suggests an odd function component plus a constant or even component. Let's try f(x) = (x\^{}2 + x)/2? No. Let's try f(x) = x*(x+1)/2? f(2)=3. No. Let's try f(x) = (x\^{}2 {-} 1)/something? Let's try f(x) = (x\^{}2 + x {-} 2)/2? f(2)=(4+2{-}2)/2=2. No. Let's try f(x) = (x\^{}2 {-} x)/2? f(2)=(4{-}2)/2=1. f({-}1)=(1+1)/2=1. But f({-}1)={-}1. So not that. Let's try f(x) = (x\^{}2 + x)/2 {-} 1? f(2)=(4+2)/2{-}1=2. No. Let's try f(x) = (x\^{}2 {-} 1)/3? f(} \\
            }
        }
    }
    & & \\ \\

    \theutterance \stepcounter{utterance}  
    & & & \multicolumn{4}{p{0.6\linewidth}}{
        \cellcolor[rgb]{0.9,0.9,0.9}{
            \makecell[{{p{\linewidth}}}]{
                \texttt{\tiny{[P1$\langle$GM]}}
                \texttt{Game Violation: ParseError: Response is not valid JSON.} \\
\texttt{Output ONLY:} \\
\texttt{One JSON with confidence\_score, current\_mode, input\_rationale, and hypothesis keys ...} \\
\texttt{OR} \\
\texttt{SOLVE: ```python ... ```} \\
            }
        }
    }
     \\ \\

    \theutterance \stepcounter{utterance}  
    & \multicolumn{4}{p{0.6\linewidth}}{
        \cellcolor[rgb]{0.561,0.737,0.561}{
            \makecell[{{p{\linewidth}}}]{
                \texttt{\tiny{[P1$\rangle$GM]}}
                \texttt{\{"current\_mode": "probing", "input\_rationale": "I have five data points: f(2)=1, f(10)=56, f({-}10)=55, f({-}1)={-}1, f(0)=0. Let's look at the differences. f(10){-}f({-}10) = 1. f(2){-}f({-}1) = 2. This suggests an odd function component plus a constant or even component. Let's try f(x) = (x\^{}2 + x)/2? No. Let's try f(x) = x*(x+1)/2? f(2)=3. No. Let's try f(x) = (x\^{}2 {-} 1)/something? Let's try f(x) = (x\^{}2 + x {-} 2)/2? f(2)=(4+2{-}2)/2=2. No. Let's try f(x) = (x\^{}2 {-} x)/2? f(2)=(4{-}2)/2=1. f({-}1)=(1+1)/2=1. But f({-}1)={-}1. So not that. Let's try f(x) = (x\^{}2 + x)/2 {-} 1? f(2)=(4+2)/2{-}1=2. No. Let's try f(x) = (x\^{}2 {-} 1)/3? f(} \\
            }
        }
    }
    & & \\ \\

    \theutterance \stepcounter{utterance}  
    & & & \multicolumn{4}{p{0.6\linewidth}}{
        \cellcolor[rgb]{0.9,0.9,0.9}{
            \makecell[{{p{\linewidth}}}]{
                \texttt{\tiny{[P1$\langle$GM]}}
                \texttt{Game Violation: ParseError: Response is not valid JSON.} \\
\texttt{Output ONLY:} \\
\texttt{One JSON with confidence\_score, current\_mode, input\_rationale, and hypothesis keys ...} \\
\texttt{OR} \\
\texttt{SOLVE: ```python ... ```} \\
            }
        }
    }
     \\ \\

    \theutterance \stepcounter{utterance}  
    & \multicolumn{4}{p{0.6\linewidth}}{
        \cellcolor[rgb]{0.561,0.737,0.561}{
            \makecell[{{p{\linewidth}}}]{
                \texttt{\tiny{[P1$\rangle$GM]}}
                \texttt{\{"current\_mode": "probing", "input\_rationale": "I have five data points: f(2)=1, f(10)=56, f({-}10)=55, f({-}1)={-}1, f(0)=0. Let's look at the differences. f(10){-}f({-}10) = 1. f(2){-}f({-}1) = 2. This suggests an odd function component plus a constant or even component. Let's try f(x) = (x\^{}2 + x)/2? No. Let's try f(x) = x*(x+1)/2? f(2)=3. No. Let's try f(x) = (x\^{}2 {-} 1)/something? Let's try f(x) = (x\^{}2 + x {-} 2)/2? f(2)=(4+2{-}2)/2=2. No. Let's try f(x) = (x\^{}2 {-} x)/2? f(2)=(4{-}2)/2=1. f({-}1)=(1+1)/2=1. But f({-}1)={-}1. So not that. Let's try f(x) = (x\^{}2 + x)/2 {-} 1? f(2)=(4+2)/2{-}1=2. No. Let's try f(x) = (x\^{}2 {-} 1)/3? f(} \\
            }
        }
    }
    & & \\ \\

    \theutterance \stepcounter{utterance}  
    & & & \multicolumn{4}{p{0.6\linewidth}}{
        \cellcolor[rgb]{0.9,0.9,0.9}{
            \makecell[{{p{\linewidth}}}]{
                \texttt{\tiny{[P1$\langle$GM]}}
                \texttt{Game Violation: ParseError: Response is not valid JSON.} \\
\texttt{Output ONLY:} \\
\texttt{One JSON with confidence\_score, current\_mode, input\_rationale, and hypothesis keys ...} \\
\texttt{OR} \\
\texttt{SOLVE: ```python ... ```} \\
            }
        }
    }
     \\ \\

    \theutterance \stepcounter{utterance}  
    & \multicolumn{4}{p{0.6\linewidth}}{
        \cellcolor[rgb]{0.561,0.737,0.561}{
            \makecell[{{p{\linewidth}}}]{
                \texttt{\tiny{[P1$\rangle$GM]}}
                \texttt{\{"current\_mode": "probing", "input\_rationale": "I have five data points: f(2)=1, f(10)=56, f({-}10)=55, f({-}1)={-}1, f(0)=0. Let's look at the differences. f(10){-}f({-}10) = 1. f(2){-}f({-}1) = 2. This suggests an odd function component plus a constant or even component. Let's try f(x) = (x\^{}2 + x)/2? No. Let's try f(x) = x*(x+1)/2? f(2)=3. No. Let's try f(x) = (x\^{}2 {-} 1)/something? Let's try f(x) = (x\^{}2 + x {-} 2)/2? f(2)=(4+2{-}2)/2=2. No. Let's try f(x) = (x\^{}2 {-} x)/2? f(2)=(4{-}2)/2=1. f({-}1)=(1+1)/2=1. But f({-}1)={-}1. So not that. Let's try f(x) = (x\^{}2 + x)/2 {-} 1? f(2)=(4+2)/2{-}1=2. No. Let's try f(x) = (x\^{}2 {-} 1)/3? f(} \\
            }
        }
    }
    & & \\ \\

    \theutterance \stepcounter{utterance}  
    & & & \multicolumn{4}{p{0.6\linewidth}}{
        \cellcolor[rgb]{0.9,0.9,0.9}{
            \makecell[{{p{\linewidth}}}]{
                \texttt{\tiny{[P1$\langle$GM]}}
                \texttt{Game Violation: ParseError: Response is not valid JSON.} \\
\texttt{Output ONLY:} \\
\texttt{One JSON with confidence\_score, current\_mode, input\_rationale, and hypothesis keys ...} \\
\texttt{OR} \\
\texttt{SOLVE: ```python ... ```} \\
            }
        }
    }
     \\ \\

    \theutterance \stepcounter{utterance}  
    & \multicolumn{4}{p{0.6\linewidth}}{
        \cellcolor[rgb]{0.561,0.737,0.561}{
            \makecell[{{p{\linewidth}}}]{
                \texttt{\tiny{[P1$\rangle$GM]}}
                \texttt{\{"current\_mode": "probing", "input\_rationale": "I have five data points: f(2)=1, f(10)=56, f({-}10)=55, f({-}1)={-}1, f(0)=0. Let's look at the differences. f(10){-}f({-}10) = 1. f(2){-}f({-}1) = 2. This suggests an odd function component plus a constant or even component. Let's try f(x) = (x\^{}2 + x)/2? No. Let's try f(x) = x*(x+1)/2? f(2)=3. No. Let's try f(x) = (x\^{}2 {-} 1)/something? Let's try f(x) = (x\^{}2 + x {-} 2)/2? f(2)=(4+2{-}2)/2=2. No. Let's try f(x) = (x\^{}2 {-} x)/2? f(2)=(4{-}2)/2=1. f({-}1)=(1+1)/2=1. But f({-}1)={-}1. So not that. Let's try f(x) = (x\^{}2 + x)/2 {-} 1? f(2)=(4+2)/2{-}1=2. No. Let's try f(x) = (x\^{}2 {-} 1)/3? f(} \\
            }
        }
    }
    & & \\ \\

    \theutterance \stepcounter{utterance}  
    & & & \multicolumn{4}{p{0.6\linewidth}}{
        \cellcolor[rgb]{0.9,0.9,0.9}{
            \makecell[{{p{\linewidth}}}]{
                \texttt{\tiny{[P1$\langle$GM]}}
                \texttt{Game Violation: ParseError: Response is not valid JSON.} \\
\texttt{Output ONLY:} \\
\texttt{One JSON with confidence\_score, current\_mode, input\_rationale, and hypothesis keys ...} \\
\texttt{OR} \\
\texttt{SOLVE: ```python ... ```} \\
            }
        }
    }
     \\ \\

    \theutterance \stepcounter{utterance}  
    & \multicolumn{4}{p{0.6\linewidth}}{
        \cellcolor[rgb]{0.561,0.737,0.561}{
            \makecell[{{p{\linewidth}}}]{
                \texttt{\tiny{[P1$\rangle$GM]}}
                \texttt{\{"current\_mode": "probing", "input\_rationale": "I have five data points: f(2)=1, f(10)=56, f({-}10)=55, f({-}1)={-}1, f(0)=0. Let's look at the differences. f(10){-}f({-}10) = 1. f(2){-}f({-}1) = 2. This suggests an odd function component plus a constant or even component. Let's try f(x) = (x\^{}2 + x)/2? No. Let's try f(x) = x*(x+1)/2? f(2)=3. No. Let's try f(x) = (x\^{}2 {-} 1)/something? Let's try f(x) = (x\^{}2 + x {-} 2)/2? f(2)=(4+2{-}2)/2=2. No. Let's try f(x) = (x\^{}2 {-} x)/2? f(2)=(4{-}2)/2=1. f({-}1)=(1+1)/2=1. But f({-}1)={-}1. So not that. Let's try f(x) = (x\^{}2 + x)/2 {-} 1? f(2)=(4+2)/2{-}1=2. No. Let's try f(x) = (x\^{}2 {-} 1)/3? f(} \\
            }
        }
    }
    & & \\ \\

    \theutterance \stepcounter{utterance}  
    & & & \multicolumn{2}{p{0.3\linewidth}}{
        \cellcolor[rgb]{0.9,0.9,0.9}{
            \makecell[{{p{\linewidth}}}]{
                \texttt{\tiny{[GM$|$GM]}}
                \texttt{Game Verdict: LOSS (turn limit reached)} \\
            }
        }
    }
    & & \\ \\

\end{supertabular}
}

\end{document}
